\begin{lstlisting}[language=Python]

def varianciaPadrao(amostras):
    media = 0  # Média dos valores
    valores = 0
    N = len(amostras)  # Quantidade de amostras

    # loop que preenche a media
    for x in amostras:
        media = media + x

    # Cálculo de média
    media = media / N

    # Iteração que efetua a soma dos valores
    for x in amostras:
        valores = valores + (x - media) ** 2

    return valores / N


# Função pra calcular media em recursão
# O índice, nesse caso, representa a quantidade de itens
def mediaWelford(amostras, indice):
    # Se o valor for o primeiro, retorna ele, pois ele ja é a media dele mesmo, se nao, continua a recursao
    if indice == 0:
        return amostras[indice]
    else:
        mediaAnterior = mediaWelford(amostras, indice - 1)
        return mediaAnterior + (amostras[indice] - mediaAnterior) / (indice + 1)


def valoresWelford(amostras, indice, media):
    # Se o valor for o primeiro, retorna o quadrado dele, se nao, continua com o calculo da recursão
    if indice == 0:
        return (amostras[indice] - media) * (amostras[indice] - media)
    else:
        return valoresWelford(amostras, indice - 1, media) + (amostras[indice] - media) * (amostras[indice] - media)


def varianciaWelford(amostras):
    # Armazena a média
    media = mediaWelford(amostras, len(amostras) - 1)

    # Armazena os valores ao quadrado somados
    somaDesviosQuadrados = valoresWelford(amostras, len(amostras) - 1, media)

    # Resultado da variância
    return somaDesviosQuadrados / len(amostras)


# Calculando...
# lista = [15, 29, 34, 40]
# print(varianciaPadrao(lista)) # Saída: 85.25
# print(varianciaWelford(lista)) # Saída: 85.25

\end{lstlisting}